\subsection{Evaluating Expressions} 
To \emph{evaluate} an expression means to determine what number it represents. When the expression contains no variables, this is straight forward. With more than one operation, we follow the following order.
\begin{enumerate}
\item Grouping symbols---including \makeblank[1.5in] and \makeblank[1.5in]---come first. These are used to indicate that some operation comes out of order.
\item Exponents are evaluated before multiplication (or division). 
\item Multiplication (and \makeblank[1.5in]) come next. Note: \makeblank[1in] signs are at this level.
\item Addition (and \makeblank[1.5in]) come last.
\end{enumerate}
For complicated expressions, this is most easily done term-by-term.
\begin{Example}
Evaluate $\frac{(7-5)^3}{4}-9\div 3\cdot 7+8$
\begin{cpic}{}{5}
\end{cpic}
\end{Example}
Negative signs can create some confusion.
\begin{itemize}
\item $-5^2$ means first \makeblank[1in], then \makeblank[1in], so $-5^2=\makeblanksmall$
\item $(-5)^2$ means first \makeblank[1in], then \makeblank[1in], so $(-5)^2=\makeblanksmall$
\end{itemize}
When an expression contains variables, we must know what \makeblank[1in] the variables represent in order to evaluate.
\begin{Example}
Evaluate $3x^2-5xy+4y^2$ when $x=-2$ and $y=3$.
\begin{lpic}{}{4}
\regtextleft{0}{-1}{First we plug in:};
\end{lpic}
\end{Example}